% !TeX root = ../动手学强化学习笔记.tex
% 由主文档通过 \input 载入；请编译主文档。
\setcounter{chapter}{12} % 保留原书章号。
\chapter{深度价值方法}
\label{ch:rl-dqn}

深度 Q 网络（deep Q-network，DQN）用神经网络表示动作价值，将 Q-learning 推广到高维观测\cite{mnih-dqn}。网络扩大了表示能力，也带来了样本相关、目标移动和自举误差传播的问题；经验回放与目标网络分别缓解其中一部分困难。

\section{从价值表到动作价值网络}

离散动作数为 $K$ 时，网络 $Q_w(s,\cdot)$ 可以一次输出 $K$ 个动作价值。交互中先按当前网络构造探索策略，环境返回转移 $(s,a,r,s',d)$，其中 $d=1$ 表示真实终止。训练仅提取动作 $a$ 对应的输出，与该转移的目标比较。若一次批量有 $B$ 条转移，提取后的预测和目标均应为长度 $B$ 的向量，避免意外广播成为 $B\times B$ 的损失矩阵。

直接把表格更新改成网络梯度不能继承表格收敛定理，因为修改参数会同时改变未采样动作的预测。连续动作下还不能简单输出并枚举所有动作，需其他优化结构。

\section{经验回放与目标网络}

经验回放（experience replay）将过去转移保存在有限容量缓冲区 $\mathcal D$，训练时从中抽取小批次。它允许一条转移被重复使用，并减弱相邻更新的时序相关性；但缓冲区仍由历史策略产生，不保证覆盖充分或全部样本独立。

在线参数为 $w$，目标参数为 $w^-$。对每条转移定义
\begin{equation}
y=r+\gamma(1-d)\max_{a'}Q_{w^-}(s',a'),\qquad
L(w)=\frac1B\sum_{i=1}^B\frac12\bigl(Q_w(s_i,a_i)-y_i\bigr)^2.
\label{eq:rl-dqn-loss}
\end{equation}
目标在此次梯度更新中停止梯度。目标参数可每隔 $C$ 次更新复制一次在线参数，也可采用 $w^-\leftarrow(1-\tau)w^-+\tau w$ 的缓慢跟踪；$C$ 或 $\tau\in(0,1]$ 必须明确。后继若只是被采样截断，$d$ 应按真实任务终止而非程序停止来赋值。

\section{双重估计与结构改进}

Double DQN 用在线网络选择后继动作，再由目标网络评价：
\[
a^*\in\arg\max_aQ_w(s',a),\qquad
y=r+\gamma(1-d)Q_{w^-}(s',a^*).
\]
普通 DQN 在同一目标网络内选择并评价动作，Double DQN 的区别正发生在这里。

优势分解网络（dueling architecture）可写为 $Q_w(s,a)=V_w(s)+A_w(s,a)-K^{-1}\sum_bA_w(s,b)$。减去动作均值消除分解中的任意平移，但此处网络内部的 $V_w$ 是动作均匀平均意义的分量，不自动等于当前策略的真实状态价值。优先经验回放按较大误差提高采样概率；若目标仍是缓冲区均匀平均损失，从概率 $p_i$ 抽样时可用 $1/(|\mathcal D|p_i)$ 校正单项梯度。实践中的部分校正、权重归一化和裁剪会改变精确对应关系。

\section{训练循环与故障诊断}

初始化网络、目标网络和回放池，先收集足够转移；每次交互后按指定频率抽取批次，计算固定目标，反向传播在线损失，再按规则同步目标网络。真实终止后重置环境，评价阶段使用预先规定的动作规则。训练预算结束时输出在线网络和评价策略，同时保存使实验可复现的配置。

若 Q 值快速增大而奖励没有增长，应检查奖励尺度、终止掩码、目标网络是否被错误反传以及后继动作维度。损失下降但回报不变可能来自覆盖不足或训练目标错误。消融应一次移除一个机制，并比较环境交互和参数更新两种成本，不能把更多梯度步数的收益全部归因于网络结构。

\begin{example}[普通与双重目标]
后继在线网络输出 $(5,4)$，目标网络输出 $(2,3)$，奖励为 $1$，$\gamma=0.9$，且未终止。普通 DQN 目标为 $1+0.9\times3=3.7$；Double DQN 先选在线最大值对应的第一动作，再评价为 $1+0.9\times2=2.8$。若真实终止，两者目标都为 $1$。
\end{example}

\section{习题与实验}
\begin{enumerate}
\item 对平方损失推导关于在线参数的梯度，说明为什么不更新目标参数，以及真实终止时后继网络输出无关。
\item 证明优势分解中各动作的平均输出等于网络分量 $V_w(s)$，并解释它与 $v_\pi(s)$ 的区别。
\item 在小离散环境先对照表格 Q-learning，再比较有无回放及目标网络的训练。报告回报、价值尺度、独立重复和更新预算。
\end{enumerate}





